\section{Cel i zakres pracy}

Celem pracy jest zaprojektowanie i zaimplementowanie aplikacji sieciowej umożliwiającej analizę biomechaniki martwego ciągu na podstawie nagrania wideo. System ma umożliwiać użytkownikowi przesłanie materiału wideo, przeprowadzenie analizy po stronie serwera oraz przedstawienie wyników w czytelnej formie za pośrednictwem interfejsu graficznego aplikacji.

Projektowane rozwiązanie wykorzystuje gotowy model estymacji pozy człowieka \cite{mediapipe_pose_landmarker}. W pracy pojęcie estymacji pozy człowieka oznacza proces wykrywania punktów kluczowych sylwetki, odpowiadających między innymi położeniom wybranych stawów wykorzystywanych w dalszej analizie \cite{mediapipe_pose_docs}. Uzyskane dane są następnie wykorzystywane do obliczenia wybranych parametrów biomechanicznych oraz oceny techniki wykonania ćwiczenia \cite{wade2022markerless,mundt2024twodimensional}.

Zakres pracy obejmuje:
\begin{itemize}
\item analizę podstawowych zagadnień związanych z biomechaniką martwego ciągu oraz oceną jego techniki,
\item omówienie możliwości wykorzystania estymacji pozy człowieka do analizy ruchu,
\item wybór oraz zastosowanie gotowego modelu wykrywania punktów kluczowych sylwetki,
\item zaprojektowanie architektury aplikacji sieciowej,
\item implementację części klienckiej odpowiedzialnej za interakcję z użytkownikiem,
\item implementację części serwerowej odpowiedzialnej za obsługę przesłanych nagrań oraz uruchamianie analizy,
\item implementację modułu analizy nagrania wideo,
\item obliczenie wybranych parametrów biomechanicznych na podstawie punktów kluczowych sylwetki,
\item opracowanie reguł oceny techniki wykonania ćwiczenia,
\item prezentację wyników analizy z podziałem na poszczególne powtórzenia oraz podsumowanie całej serii,
\item przeprowadzenie testów działania systemu na wybranych nagraniach.
\end{itemize}

Praca ma charakter aplikacyjny. Jej celem nie jest opracowanie nowego modelu uczenia maszynowego ani trenowanie modelu od podstaw. Główny nacisk położono na wykorzystanie istniejącego modelu estymacji pozy człowieka oraz przygotowanie kompletnego systemu informatycznego integrującego moduł analizy nagrania wideo z aplikacją sieciową.

W pracy przyjęto również kilka ograniczeń projektowych. System został przygotowany pod kątem analizy nagrań spełniających następujące warunki:
\begin{itemize}
\item nagranie wykonano z perspektywy bocznej,
\item kamera pozostaje statyczna,
\item przez cały czas widoczna jest cała sylwetka osoby wykonującej ćwiczenie,
\item w kadrze widoczna jest tylko jedna osoba,
\item rozmiar pliku oraz czas trwania nagrania mieszczą się w przyjętych ograniczeniach,
\item nagranie zapisano w obsługiwanym formacie pliku.
\end{itemize}

Takie wymagania ułatwiają analizę ruchu w płaszczyźnie strzałkowej, w szczególności analizę ustawienia tułowia oraz pracy stawów biodrowych i kolanowych \cite{hanen2025deadlift,wade2022markerless}. Jakość materiału wejściowego ma istotny wpływ na uzyskiwane wyniki, co zostanie dokładniej omówione w dalszej części pracy \cite{wade2022markerless}.

Projektowany system należy traktować jako narzędzie wspomagające analizę techniki, a nie jako system diagnostyczny gwarantujący całkowitą dokładność oceny. Wyniki generowane przez aplikację mają pomóc w ocenie wybranych aspektów ruchu i wskazać potencjalne błędy, jednak nie zastępują profesjonalnej analizy prowadzonej przez trenera, fizjoterapeutę lub specjalistę z zakresu biomechaniki \cite{wade2022markerless}.
